\chapter{Espaços de Banach e os teoremas fundamentais}
\label{ch:b3:banach}

A análise funcional estuda os espaços normados de dimensão infinita por
meio dos operadores e funcionais que vivem sobre eles. Sua descoberta
fundadora é que a \emph{\omterm{def:b3:complete:complete}{completude}}, via o teorema de Baire, impõe
fortes uniformidades: famílias de operadores pontualmente limitadas são
limitadas em norma (Banach--Steinhaus), bijeções \omterm{def:b3:topology:continuity}{contínuas} têm inversas
\omterm{def:b3:topology:continuity}{contínuas} (aplicação aberta) e os gráficos detectam a \omterm{def:b3:topology:continuity}{continuidade}
(gráfico fechado). O outro pilar, \emph{Hahn--Banach}, não exige
\omterm{def:b3:complete:complete}{completude} alguma --- apenas o lema de Zorn --- e garante que os
\omterm{def:b3:banach:operator}{espaços duais} sejam ricos o bastante para enxergar todo vetor. Este
capítulo demonstra os quatro teoremas e os testa nos espaços clássicos
de sequências $\ell^p$, em cálculos concretos de duais e em uma
aplicação genuinamente surpreendente: existem funções \omterm{def:b3:topology:continuity}{contínuas}
$2\pi$-periódicas cuja série de Fourier \emph{diverge} em um ponto
--- resolvendo negativamente uma questão que o segundo ano deixou em
aberto.

Ao longo de todo o capítulo, $E, F$ são espaços normados sobre
$K = \R$ ou $\C$; \emph{espaço de Banach}\index{espaço de Banach}
significa espaço normado \omterm{def:b3:complete:complete}{completo}.

\section{Operadores limitados; espaços de sequências}

\begin{definition}\label{def:b3:banach:operator}
$\mathcal L(E, F)$ denota o espaço das aplicações lineares \emph{limitadas}
(= \omterm{def:b3:topology:continuity}{contínuas}, segundo ano) com a \emph{norma de operador}\index{norma
de operador} $\vertiii T = \sup_{\norm x \leq
1}\norm{Tx}$; ela é submultiplicativa: $\vertiii{ST} \leq
\vertiii S\,\vertiii T$. O
\emph{dual}\index{espaço dual} é $E' = \mathcal L(E, K)$.
\end{definition}

\begin{proposition}\label{prop:b3:banach:llcomplete}
Se $F$ é um espaço de Banach, $\mathcal L(E, F)$ também é; em particular, $E'$ é
sempre um espaço de Banach.
\end{proposition}

\begin{proof}
Seja $(T_n)$ de Cauchy para $\vertiii\cdot$. Para cada $x$, $\norm{T_nx - T_mx} \leq \vertiii{T_n - T_m}\norm x$: $(T_nx)$
é de Cauchy em $F$ e, portanto, convergente; chame o limite de
$Tx$. $T$ é linear (limites de identidades lineares); passando ao
limite em $\norm{T_nx - T_mx} \leq \varepsilon\norm x$ ($n, m \geq N$) obtém-se $\norm{T_nx - Tx} \leq \varepsilon\norm x$: $T_n \to T$ em \omterm{def:b3:banach:operator}{norma de operador}, e $\vertiii T \leq \vertiii{T_N} +
\varepsilon < \infty$.
\end{proof}

\begin{definition}\label{def:b3:banach:lp}
Os espaços clássicos de sequências\index{espaços de sequências
$\ell^p$} (sobre $K$, indexados por $\N$):
\[
\ell^p = \Bigl\{x = (x_n) : \norm x_p = \Bigl(\sum_n
\abs{x_n}^p\Bigr)^{1/p} < \infty\Bigr\}\quad (1 \leq p <
\infty),
\]
\[
\ell^\infty = \{x : \norm x_\infty = \sup\abs{x_n} < \infty\},
\]
e $c_0 = \{x : x_n \to 0\}$ com $\norm\cdot_\infty$. Que $\norm\cdot_p$ é uma norma decorre da desigualdade de
Minkowski, demonstrada no caso discreto exatamente como no~\cref{ch:b3:lp} (ou somando a desigualdade em dimensão finita do segundo
ano). Todos são espaços de Banach, e $c_0$ é um subespaço fechado de
$\ell^\infty$ (\cref{exo:b3:banach:3}).
\end{definition}

\begin{omfigure}
\begin{tikzpicture}[scale=1.5]
  \draw[->] (-1.45,0) -- (1.45,0) node[right,font=\small]
    {$x_1$};
  \draw[->] (0,-1.45) -- (0,1.45) node[above,font=\small]
    {$x_2$};
  % p = 1 diamond
  \draw[omDef, thick] (1,0) -- (0,1) -- (-1,0) -- (0,-1)
    -- cycle;
  % p = 2 circle
  \draw[omProp, thick] (0,0) circle (1);
  % p = 4 superellipse |x|^4 + |y|^4 = 1
  \draw[omThm, thick, smooth]
    plot[domain=0:90, samples=46]
    ({cos(\x)/(cos(\x)^4 + sin(\x)^4)^0.25},
     {sin(\x)/(cos(\x)^4 + sin(\x)^4)^0.25});
  \draw[omThm, thick, smooth]
    plot[domain=90:180, samples=46]
    ({cos(\x)/(cos(\x)^4 + sin(\x)^4)^0.25},
     {sin(\x)/(cos(\x)^4 + sin(\x)^4)^0.25});
  \draw[omThm, thick, smooth]
    plot[domain=180:270, samples=46]
    ({cos(\x)/(cos(\x)^4 + sin(\x)^4)^0.25},
     {sin(\x)/(cos(\x)^4 + sin(\x)^4)^0.25});
  \draw[omThm, thick, smooth]
    plot[domain=270:360, samples=46]
    ({cos(\x)/(cos(\x)^4 + sin(\x)^4)^0.25},
     {sin(\x)/(cos(\x)^4 + sin(\x)^4)^0.25});
  % p = infty square
  \draw[black!60, thick] (-1,-1) rectangle (1,1);
  \node[omDef, font=\small] at (0.40,0.40) {$p{=}1$};
  \node[omProp, font=\small] at (-0.52,0.75) {$p{=}2$};
  \node[omThm, font=\small] at (0.90,-0.52) {$p{=}4$};
  \node[black!60, font=\small] at (-1.62,1.15)
    {$p{=}\infty$};
  \draw[black!60, ->] (-1.32,1.11) -- (-1.02,1.0);
\end{tikzpicture}

{\small Bolas unitárias das normas $p$ no plano, encaixadas à medida
que $p$ cresce de $1$ (losango), passando por $2$ (disco) e
$4$ (superelipse), até $\infty$ (quadrado). A convexidade de cada
bola \emph{é} a desigualdade de Minkowski; os cantos em $p = 1$ e
$p = \infty$ são onde a convexidade estrita, a unicidade das melhores
aproximações e os casos de igualdade do~\cref{exo:b3:lp:12} degeneram todos de
uma só vez.}
\end{omfigure}

\begin{proposition}[Série de Neumann]\label{prop:b3:banach:neumann}
Sejam $E$ de Banach e $T \in \mathcal L(E) = \mathcal L(E,E)$ com $\vertiii T < 1$. Então $I - T$ é inversível em
$\mathcal
L(E)$, com $(I - T)^{-1} = \sum_{n\geq0}T^n$ (convergente em \omterm{def:b3:banach:operator}{norma de operador}). Consequentemente, o
conjunto dos operadores inversíveis é aberto, e a inversão é \omterm{def:b3:topology:continuity}{contínua}
nele.\index{série de Neumann}
\end{proposition}

\begin{proof}
A série converge absolutamente ($\vertiii{T^n} \leq \vertiii
T^n$, geométrica) no espaço de Banach
$\mathcal L(E)$ (\cref{prop:b3:banach:llcomplete}; \cref{exo:b3:complete:1}(b)). Telescopando, $(I - T)\sum_{n \leq
N}T^n = I - T^{N+1} \to I$, e
analogamente do outro lado. Quanto à abertura: se $S$ é inversível e
$\vertiii{H} <
1/\vertiii{S^{-1}}$, então $S + H = S(I + S^{-1}H)$ com $\vertiii{S^{-1}H} < 1$: inversível. \omterm{def:b3:topology:continuity}{Continuidade} da inversão: a
expressão em série dá $\vertiii{(S+H)^{-1} - S^{-1}} =
O(\vertiii H)$ localmente.
\end{proof}

\begin{example}[Uma equação de Volterra, resolvida por Neumann]
\label{ex:b3:banach:volterraexample}
Em $E = \mathcal C(\intcc01)$, considere a equação integral
\[
u(x) = 1 + \lambda\int_0^xu(t)\,\dd t,
\qquad\text{isto é}\qquad u = \mathbf 1 + \lambda Tu,
\quad (Tu)(x) = \int_0^xu .
\]
Aqui $\vertiii T \leq 1$, de modo que, para $\abs\lambda < 1$, a \omterm{prop:b3:banach:neumann}{série de Neumann} se aplica
diretamente: $u = (I - \lambda T)^{-1}\mathbf 1 =
\sum_n\lambda^nT^n\mathbf 1$. Calculando, $T^n\mathbf 1 =
\frac{x^n}{n!}$, logo
\[
u(x) = \sum_{n\geq0}\frac{(\lambda x)^n}{n!} =
\eu^{\lambda x} ,
\]
como a derivação confirma. Melhor ainda: $\vertiii{T^n} \leq
\frac1{n!}$ (o núcleo iterado encolhe
fatorialmente), de modo que $\sum\lambda^nT^n$ converge para \emph{todo} $\lambda$ --- o
operador $I - \lambda T$ é inversível para todo $\lambda
\in \C$, ainda que $\vertiii{\lambda T} \geq 1$ acabe
ocorrendo: o que importa é o decaimento espectral das potências, não a
primeira norma. Os operadores de Volterra trazem esse decaimento
fatorial embutido (\cref{exo:b3:complete:4}(a) explorou exatamente isso), e é por
isso que os problemas de valor inicial nunca sofrem os fenômenos de
ressonância dos problemas de valores de contorno (\cref{ch:b3:spectral}).
\end{example}

\section{Hahn--Banach}

\begin{theorem}[Hahn--Banach, forma analítica]
\label{thm:b3:banach:hahnbanach}
Sejam $E$ um espaço vetorial \emph{real}, $p \colon E \to \R$ \emph{sublinear}
($p(x + y) \leq p(x) + p(y)$ e $p(tx) =
tp(x)$ para $t \geq 0$), $F \subseteq E$ um subespaço e $f
\colon F \to \R$ linear com $f \leq p$ em
$F$. Então $f$ se estende a uma aplicação linear $\tilde f \colon E \to \R$ com $\tilde f \leq p$ em
$E$.\index{teorema de Hahn--Banach}
\end{theorem}

\begin{proof}
\emph{Extensão de um passo.} Seja $x_0 \notin F$; estendemos $f$ a $F \oplus \R x_0$
escolhendo $\alpha = \tilde f(x_0)$ corretamente: precisamos que, para todos $y \in F$ e
$t > 0$,
\[
f(y) + t\alpha \leq p(y + tx_0)
\quad\text{e}\quad
f(y) - t\alpha \leq p(y - tx_0),
\]
o que, após dividir por $t$ (sublinearidade), se reduz a
\[
\sup_{v \in F}\ \bigl[f(v) - p(v - x_0)\bigr]
\;\leq\; \alpha \;\leq\;
\inf_{u \in F}\ \bigl[p(u + x_0) - f(u)\bigr].
\]
Um tal $\alpha$ existe se, e somente se, todo membro esquerdo é
$\leq$ a todo membro direito: com efeito, $f(v) + f(u) = f(u + v) \leq p(u + v) \leq
p(u + x_0) + p(v - x_0)$, isto é, $f(v) - p(v - x_0) \leq p(u +
x_0) - f(u)$.

\emph{Zorn.} Ordene por extensão as extensões de $f$ dominadas por
$p$ (pares: subespaço, funcional); uma cadeia tem a reunião por cota
superior; um elemento maximal tem de estar definido em todo $E$,
pois do contrário a extensão de um passo contradiria a maximalidade.
\end{proof}

\begin{corollary}\label{cor:b3:banach:hbnormed}
Seja $E$ um espaço normado ($K = \R$ ou $\C$).
\begin{enumerate}
  \item Todo $f \in F'$ (com $F$ subespaço) se estende a $\tilde f
        \in E'$ com $\norm{\tilde f}_{E'} = \norm f_{F'}$.
  \item Para todo $x \neq 0$ existe $f \in E'$ com $\norm f
        = 1$ e $f(x) = \norm x$. Em particular,
        $E'$ separa os pontos de $E$, e $\norm x = \sup_{\norm f \leq
        1}\abs{f(x)}$.
  \item Para um subespaço fechado $F$ e $x \notin F$, existe $f
        \in E'$ que se
        anula em $F$, com $f(x) = d(x, F)$ e $\norm f \leq 1$.
\end{enumerate}
\end{corollary}

\begin{proof}
(1) Caso real: aplique o~\cref{thm:b3:banach:hahnbanach} com $p(x)
= \norm f_{F'}\,\norm x$ (sublinear); a
extensão satisfaz $\pm\tilde f(x) = \tilde f(\pm x) \leq p(x)$, logo $\norm{\tilde f} \leq \norm f$, e $\geq$ vem da restrição. Caso
complexo: ponha $u = \operatorname{Re}f$, um funcional real com $\abs u \leq \norm f\norm\cdot$; note que $f(x) = u(x) -
\iu\,u(\iu x)$
(verifique nas partes real e imaginária: $\operatorname{Im}f(x) = -\operatorname{Re}f(\iu x)$). Estenda $u$ de
maneira real-linear com a mesma cota e ponha $\tilde f(x) =
\tilde u(x) - \iu\tilde u(\iu x)$: é $\C$-linear
(verificação direta na multiplicação por $\iu$) e estende $f$;
quanto à norma, para $x$ dado escreva $\tilde f(x) = r\eu^{\iu\theta}$, e então $\abs{\tilde f(x)}
= \tilde f(\eu^{-\iu\theta}x) = \tilde u(\eu^{-\iu\theta}x)
\leq \norm f\,\norm x$.

(2) Em $F = Kx$, defina $f(tx) = t\norm x$: de norma $1$ em $F$; estenda
por (1). A fórmula de dualidade: $\leq$ é claro, e $\geq$ vem
desse $f$.

(3) Em $F \oplus Kx$, defina $f(y + tx) = t\,d(x, F)$; então, para $t \neq 0$, $\norm{y + tx} = \abs t\,\norm{x + y/t} \geq \abs
t\,d(x, F) = \abs{f(y + tx)}$: $\norm f \leq 1$ no
subespaço; estenda por (1).
\end{proof}

\begin{remark}\label{rem:b3:banach:bidual}
Por (2), a aplicação canônica $J \colon E \to E''$, $J(x)(f) =
f(x)$, é uma isometria
(\cref{exo:b3:banach:10}): todo espaço normado se aloja dentro de seu bidual. Os
espaços com $J$ sobrejetora são chamados
\emph{reflexivos}\index{espaço reflexivo}; o problema de fim de semana
mostra que $\ell^p$ ($1 < p < \infty$) é \omterm{rem:b3:banach:bidual}{reflexivo}, ao passo que $\ell^1$
não é.
\end{remark}

\section{A trilogia de Baire}

\begin{theorem}[Banach--Steinhaus, limitação uniforme]
\label{thm:b3:banach:banachsteinhaus}
Sejam $E$ um espaço \emph{de Banach}, $F$ normado e $(T_i)_{i\in
I} \subseteq \mathcal L(E, F)$ uma
família com $\sup_i
\norm{T_ix} < \infty$ para todo $x \in E$. Então $\sup_i
\vertiii{T_i} < \infty$.\index{teorema de
Banach--Steinhaus}
\end{theorem}

\begin{proof}
Os conjuntos $F_n = \{x : \sup_i\norm{T_ix} \leq n\}$ são fechados (interseções de pré-imagens de bolas
fechadas) e cobrem $E$. Baire (\cref{thm:b3:complete:baire}) dá $n_0$ e uma
bola $B(x_0, r) \subseteq F_{n_0}$. Para $\norm z < r$: $\norm{T_iz}
\leq \norm{T_i(x_0 + z)} + \norm{T_ix_0} \leq 2n_0$, logo $\vertiii{T_i} \leq 2n_0/r$ para todo $i$.
\end{proof}

\begin{corollary}\label{cor:b3:banach:bslimits}
Se $E$ é de Banach e $T_n \in \mathcal L(E,F)$ convergem \emph{pontualmente} ($T_nx \to Tx$ para
cada $x$), então $\sup_n
\vertiii{T_n} < \infty$, $T \in \mathcal L(E, F)$, e $\vertiii T \leq \liminf \vertiii{T_n}$.
\end{corollary}

\begin{proof}
Sequências convergentes são limitadas: eis a limitação pontual;
Banach--Steinhaus limita as normas por algum $M$; então $\norm{Tx}
= \lim\norm{T_nx} \leq M\norm x$
($T$ é linear por ser limite pontual), e a cota mais fina segue
passando a $\liminf$ em $\norm{T_nx} \leq \vertiii{T_n}\norm x$.
\end{proof}

\begin{theorem}[Séries de Fourier divergentes]
\label{thm:b3:banach:fourierdiverge}
Existem funções \omterm{def:b3:topology:continuity}{contínuas} $2\pi$-periódicas $f$ cuja série de
Fourier diverge em $0$: $\sup_N\abs{S_N(f)(0)} =
\infty$. Mais ainda, tais $f$ formam um
subconjunto denso de $\mathcal
C(S^1)$.\index{série de Fourier, divergência}
\end{theorem}

\begin{proof}
Trabalhamos em $E = (\mathcal C(S^1), \norm\cdot_\infty)$, um espaço de Banach, com os funcionais $\Lambda_N(f) = S_N(f)(0) =
\frac1{2\pi}\int_{-\pi}^{\pi}f(t)\,D_N(t)\,\dd t$
(núcleo de Dirichlet, segundo ano). Cada $\Lambda_N$ é \omterm{def:b3:topology:continuity}{contínuo}, com
\[
\norm{\Lambda_N} = \frac1{2\pi}\int_{-\pi}^\pi\abs{D_N(t)}\,\dd
t \;=\; L_N .
\]
($\leq$ é claro; $\geq$: tome $f$ \omterm{def:b3:topology:continuity}{contínua}, $\norm f_\infty
\leq 1$,
aproximando $\operatorname{sign}D_N$ --- o sinal tem um número finito de saltos; suavizar
cada salto em um intervalo de comprimento $\varepsilon$ altera a integral de
$O(N\varepsilon)$.) As \emph{constantes de Lebesgue} $L_N$ tendem ao infinito:
\[
L_N = \frac1{2\pi}\int_{-\pi}^{\pi}
\frac{\abs{\sin\bigl((N{+}\tfrac12)t\bigr)}}{\abs{\sin(t/2)}}
\,\dd t
\geq \frac{2}{\pi}\int_0^{\pi}
\frac{\abs{\sin\bigl((N{+}\tfrac12)t\bigr)}}{t}\,\dd t
= \frac{2}{\pi}\int_0^{(N+\frac12)\pi}\frac{\abs{\sin
u}}{u}\,\dd u
\]
(usando $\abs{\sin(t/2)} \leq t/2$ em $[0, \pi]$ e substituindo em seguida $u = (N + \tfrac12)t$). Cortando em
arcos:
\[
\int_{(k-1)\pi}^{k\pi}\frac{\abs{\sin u}}u\,\dd u
\geq \frac1{k\pi}\int_{(k-1)\pi}^{k\pi}\abs{\sin u}\,\dd u =
\frac{2}{k\pi},
\qquad\text{logo}\qquad
L_N \geq \frac{4}{\pi^2}\sum_{k=1}^N\frac1k
\xrightarrow[N\to\infty]{} \infty .
\]
Se toda $f$ \omterm{def:b3:topology:continuity}{contínua} tivesse $\sup_N
\abs{\Lambda_N(f)} < \infty$, Banach--Steinhaus forçaria
$\sup_N\norm{\Lambda_N} < \infty$: contradição. Logo alguma $f$ --- de fato, um conjunto não
\omterm{rem:b3:complete:meagre}{magro} e denso de $f$ (o complementar de $\bigcup_M\{f: \sup_N\abs{\Lambda_Nf}\leq M\}$, uma reunião
enumerável de fechados que, não tendo \omterm{def:b3:topology:interior}{interior} pelo que precede
aplicado em qualquer bola, é magra) --- satisfaz $\sup_N\abs{S_Nf(0)} =
\infty$.
\end{proof}

\begin{theorem}[Aplicação aberta]\label{thm:b3:banach:openmapping}
Sejam $E, F$ espaços de Banach e $T \in \mathcal L(E, F)$ \emph{sobrejetora}. Então
$T$ é aberta: $T(B_E(0,1)) \supseteq
B_F(0, c)$ para algum $c > 0$. Consequentemente, um
operador limitado bijetor entre espaços de Banach tem inversa
limitada.\index{teorema da aplicação aberta}
\end{theorem}

\begin{proof}
Escreva $B = B_E(0,1)$. A sobrejetividade dá $F = \bigcup_n
\overline{T(nB)} = \bigcup_n n\,\overline{T(B)}$; Baire (\cref{thm:b3:complete:baire}) dá
\omterm{def:b3:topology:interior}{interior} a $\overline{T(B)}$: algum $B_F(y_0, 4c) \subseteq
\overline{T(B)}$. Recentrando em $0$: para $\norm y < 4c$, tanto
$y_0 + y$ quanto $y_0$ são limites de imagens $Tu_k$, $Tv_k$ com
$u_k, v_k \in B$, de modo que $y = \lim T(u_k - v_k)$ com $u_k - v_k
\in 2B$: $B_F(0, 4c) \subseteq \overline{T(2B)}$, isto é, $B_F(0, 2c) \subseteq \overline{T(B)}$.

\emph{Removendo o \omterm{def:b3:topology:interior}{fecho}} (é aqui que entra a \omterm{def:b3:complete:complete}{completude} de $E$):
seja $\norm y < c$. Tome $x_1 \in \frac12 B$ com $\norm{y -
Tx_1} < c/2$ ($B_F(0,2c) \subseteq \overline{T(B)}$ reescalado por $\frac12$); e,
indutivamente, $x_k \in 2^{-k}B$ com $\norm{y -
T(x_1 + \dots + x_k)} < c\,2^{-k}$. A série $\sum x_k$ converge absolutamente no
espaço de Banach $E$, para um $x \in B$ (de norma $<
\sum 2^{-k} = 1$), e $Tx = y$
por \omterm{def:b3:topology:continuity}{continuidade}: $B_F(0, c)
\subseteq T(B)$. A abertura de $T$ em abertos arbitrários
segue por translação e reescalamento; quanto ao corolário, ser $T$
aberta significa que $T^{-1}$ é \omterm{def:b3:topology:continuity}{contínua}.
\end{proof}

\begin{corollary}[Normas equivalentes]\label{cor:b3:banach:equivnorms}
Se um espaço vetorial é \omterm{def:b3:complete:complete}{completo} para duas normas comparáveis ($\norm\cdot_a \leq C\norm\cdot_b$),
então as normas são equivalentes.
\end{corollary}

\begin{proof}
A identidade $(E, \norm\cdot_b) \to (E, \norm\cdot_a)$ é limitada e bijetora entre espaços de Banach: sua
inversa é limitada.
\end{proof}

\begin{theorem}[Gráfico fechado]\label{thm:b3:banach:closedgraph}
Sejam $E, F$ de Banach e $T \colon E \to F$ linear. Se o gráfico $\Gamma = \{(x, Tx)\}$ é fechado
em $E \times F$ (isto é, $x_n \to x$ e $Tx_n \to y$ implicam $y = Tx$), então $T$ é
limitada.\index{teorema do gráfico fechado}
\end{theorem}

\begin{proof}
$E \times F$ com $\norm{(x,y)} = \norm x + \norm y$ é de Banach; $\Gamma$, subespaço fechado, é de
Banach. A projeção $\pi_E\colon \Gamma \to E$ é limitada e bijetora, de modo que sua inversa
$x \mapsto (x, Tx)$ é limitada (\cref{thm:b3:banach:openmapping}): $\norm{Tx} \leq \norm{(x,
Tx)} \leq C\norm x$.
\end{proof}

\begin{method}\label{met:b3:banach:trilogy}
Quando recorrer a qual teorema. \emph{Hahn--Banach}: para produzir um
funcional com comportamento prescrito (normar um vetor, anular-se em um
subespaço, estender de um subespaço) --- não é preciso \omterm{def:b3:complete:complete}{completude}
alguma. \emph{Banach--Steinhaus}: para converter informação pontual em
cotas uniformes --- tipicamente, para mostrar que uma operação de
limite é \omterm{def:b3:topology:continuity}{contínua} ou, na contrapositiva, para demonstrar
\emph{divergência} para algum elemento, como no caso das séries de
Fourier. \emph{Aplicação aberta / gráfico fechado}: para obter
\omterm{def:b3:topology:continuity}{continuidade} de graça a partir da bijetividade \omterm{def:b3:galois:algebraic}{algébrica} ou de uma
propriedade de \omterm{def:b3:topology:interior}{fecho} do gráfico --- uso típico: comparar duas normas
\omterm{def:b3:complete:complete}{completas}, ou demonstrar \omterm{def:b3:topology:continuity}{continuidade} automática. Os três teoremas de
Baire exigem \omterm{def:b3:complete:complete}{completude} \emph{da fonte}; do contrário há contraexemplos
(\cref{exo:b3:banach:7}).
\end{method}

\section{Espaços duais, concretamente}

\begin{theorem}\label{thm:b3:banach:duals}
Isometricamente: $(c_0)' \cong \ell^1$ e $(\ell^1)' \cong
\ell^\infty$, por meio do emparelhamento
$\langle x, y\rangle = \sum_n
x_ny_n$.\index{duais dos espaços de sequências}
\end{theorem}

\begin{proof}
Demonstramos $(c_0)' \cong \ell^1$; a segunda identificação é o~\cref{exo:b3:banach:5}. A $y \in \ell^1$
associe $\Lambda_y(x)
= \sum x_ny_n$ ($x \in c_0$): absolutamente convergente, com $\abs{\Lambda_y(x)} \leq \norm x_\infty\norm y_1$, de modo que
$\norm{\Lambda_y} \leq \norm y_1$. Reciprocamente, seja $\Lambda \in
(c_0)'$; ponha $y_n = \Lambda(e_n)$ (com $e_n$ as
sequências unitárias). Para todo $N$, teste em $x^{(N)} = \sum_{n \leq N}
\operatorname{sign}(\overline{y_n})\,e_n \in c_0$ (de norma
$\leq
1$; no caso complexo, use fatores unimodulares $\bar y_n/\abs
{y_n}$): $\Lambda(x^{(N)}) = \sum_{n\leq N}\abs{y_n} \leq
\norm\Lambda$.
Logo $y \in \ell^1$ com $\norm y_1 \leq
\norm\Lambda$. Por fim, $\Lambda = \Lambda_y$: ambos coincidem nos $e_n$,
portanto nas sequências finitas, densas em $c_0$ (truncamento:
$\norm{x - \sum_{n \leq N}x_ne_n}_\infty = \sup_{n > N}\abs{x_n}
\to 0$ precisamente porque $x_n \to 0$); e funcionais \omterm{def:b3:topology:continuity}{contínuos} que coincidem
em um \omterm{def:b3:topology:interior}{conjunto denso} são iguais. A correspondência é linear, bijetora e
isométrica ($\norm{\Lambda_y} = \norm y_1$, pelas duas desigualdades).
\end{proof}

\section{Exercícios}

\begin{exercise}[$\star$]\label{exo:b3:banach:1}
Calcule as \omterm{def:b3:banach:operator}{normas de operador}:
(a) dos deslocamentos $S(x_1, x_2, \dots) = (0, x_1, x_2, \dots)$ e $S^*(x_1, x_2, \dots) = (x_2, x_3, \dots)$ em $\ell^2$;
(b) do operador de multiplicação $M_a x = (a_nx_n)$ em $\ell^2$, para $a \in \ell^\infty$;
(c) do funcional $\Lambda(f) = \int_0^{1/2}f -
\int_{1/2}^1f$ em $\mathcal C(\intcc01)$ --- mostre que $\norm\Lambda = 1$ e que a norma
\emph{não é atingida}.
\end{exercise}

\begin{exercise}[$\star$]\label{exo:b3:banach:2}
Sejam $E$ de Banach, $T \in \mathcal L(E)$ inversível e $S$ com $\vertiii{S - T} < 1/\vertiii{T^{-1}}$. Mostre que
$S$ é inversível e estime $\vertiii{S^{-1} - T^{-1}}$. Aplicação: se um sistema linear
$Tx = b$ é \omterm{def:b3:groups:derived}{solúvel} com $T$ inversível, uma perturbação
suficientemente pequena de $T$ o mantém unicamente \omterm{def:b3:groups:derived}{solúvel}, com uma
cota quantitativa para a variação da solução.
\end{exercise}

\begin{exercise}[$\star\star$]\label{exo:b3:banach:3}
(a) Demonstre que $\ell^1$, $\ell^\infty$ e $c_0$ são espaços de Banach,
e que $c_0$ é o \omterm{def:b3:topology:interior}{fecho} em $\ell^\infty$ do espaço das sequências finitas.
(b) Mostre que $\ell^p \subseteq \ell^q$ com $\norm\cdot_q \leq
\norm\cdot_p$ para $1 \leq p \leq q \leq \infty$, e que a inclusão é estrita.
\end{exercise}

\begin{exercise}[$\star\star$]\label{exo:b3:banach:4}
Sejam $F \subseteq E$ um subespaço fechado e $x \notin F$. Usando o~\cref{cor:b3:banach:hbnormed},
demonstre a fórmula de dualidade
\[
d(x, F) = \max\bigl\{\abs{f(x)} : f \in E',\ \norm f \leq 1,\
f\restriction_F = 0\bigr\}
\]
(atenção: um \emph{máximo}). Deduza que $F = \bigcap\{\ker f : f
\in E',\ f\restriction_F = 0\}$: os subespaços fechados
são exatamente as interseções de núcleos de funcionais.
\end{exercise}

\begin{exercise}[$\star\star$]\label{exo:b3:banach:5}
Demonstre $(\ell^1)' \cong \ell^\infty$ isometricamente, seguindo o esquema do~\cref{thm:b3:banach:duals}
(as sequências finitas são densas em $\ell^1$). Onde o argumento se
quebra para $(\ell^\infty)'$?
\end{exercise}

\begin{exercise}[$\star\star$]\label{exo:b3:banach:6}
Sejam $E, F, G$ normados com $E$ de Banach, e $B \colon E \times
F \to G$ bilinear, \omterm{def:b3:topology:continuity}{contínua}
em cada variável separadamente. Mostre que $B$ é (conjuntamente)
\omterm{def:b3:topology:continuity}{contínua}: $\norm{B(x,y)} \leq
C\norm x\norm y$. \emph{(Aplique Banach--Steinhaus à família $(B(\cdot, y))_{\norm y \leq 1}$.)}
\end{exercise}

\begin{exercise}[$\star\star$]\label{exo:b3:banach:7}
(a) Em $E = \mathcal C(\intcc01)$, compare $\norm\cdot_\infty$ e $\norm\cdot_1$: a identidade $(E, \norm\cdot_\infty) \to
(E, \norm\cdot_1)$ é limitada e
bijetora, mas sua inversa é ilimitada. Que hipótese do~\cref{cor:b3:banach:equivnorms}
falha?
(b) Exiba uma aplicação linear descontínua de um subespaço denso de
$\ell^2$ em $K$ (por exemplo, nas sequências finitas) e explique
por que isso não contradiz o teorema do gráfico fechado.
\end{exercise}

\begin{exercise}[$\star\star\star$]\label{exo:b3:banach:8}
(Hellinger--Toeplitz) Seja $T \colon \ell^2 \to \ell^2$ linear (definida em toda parte) e
\emph{simétrica}: $\langle Tx,
y\rangle = \langle x, Ty\rangle$ para todos $x, y$, em que $\langle
x, y \rangle = \sum x_n\bar y_n$. Mostre que
$T$ é limitada. \emph{(Gráfico fechado: se $x_k \to x$ e $Tx_k \to z$, teste
contra $y$ arbitrário.)} Moral: operadores simétricos ilimitados ---
as hamiltonianas da mecânica quântica --- jamais podem ser definidos no
espaço todo.
\end{exercise}

\begin{exercise}[$\star\star\star$]\label{exo:b3:banach:9}
(Teorema de Pólya sobre quadraturas) Para cada $n$, seja $\Lambda_n(f)
= \sum_{i=0}^{n} w_{i,n}f(x_{i,n})$ uma
regra de quadratura em $\mathcal C(\intcc01)$ ($x_{i,n} \in \intcc01$, $w_{i,n} \in
\R$). Mostre que $\Lambda_n(f) \to \int_0^1f$ para
\emph{toda} $f$ \omterm{def:b3:topology:continuity}{contínua} se, e somente se: (i) $\Lambda_n(P) \to \int_0^1P$ para todo
polinômio $P$, e (ii) $\sup_n\sum_i\abs{w_{i,n}} <
\infty$. \emph{(Calcule $\norm{\Lambda_n}$; use
Banach--Steinhaus e Weierstrass.)} Verifique que as regras com pesos
\emph{positivos} exatas nas constantes satisfazem (ii)
automaticamente.
\end{exercise}

\begin{exercise}[$\star\star$]\label{exo:b3:banach:10}
Mostre que $J \colon E \to E''$, $J(x)(f) = f(x)$, é uma isometria linear (use
o~\cref{cor:b3:banach:hbnormed}(2)) e que ela é sobrejetora quando $\dim E < \infty$. Mostre também
que, se $E'$ é separável, então $E$ também é. \emph{(Tome
$x_n$ que quase normem uma sequência densa de $E'$ e mostre que o
subespaço fechado que geram é $E$, via o~\cref{cor:b3:banach:hbnormed}(3).)}
\end{exercise}

\begin{exercise}[$\star\star$]\label{exo:b3:banach:11}
(Espaços quocientes) Sejam $E$ um espaço de Banach e $F \subseteq
E$ um
subespaço \emph{fechado}. Em $E/F$, defina
\[
\norm{\bar x} \;=\; d(x, F) = \inf_{y\in F}\norm{x - y} .
\]
(a) Mostre que isso define uma \emph{norma} em $E/F$ (onde entra o
fato de $F$ ser fechado?) e que a projeção $\pi
\colon E \to E/F$ satisfaz $\vertiii\pi \leq 1$ e
leva a bola unitária aberta sobre a bola unitária aberta.
(b) Mostre que $E/F$ é \omterm{def:b3:complete:complete}{completo}. \emph{(Use o critério de séries
do~\cref{exo:b3:complete:1}(b): dadas classes $\bar x_k$ com $\sum\norm{\bar x_k} < \infty$, levante cada uma a
$x_k \in E$ com $\norm{x_k} \leq \norm{\bar x_k} + 2^{-k}$ e some em $E$.)}
(c) Calcule: para $E = c$ (sequências convergentes) e $F =
c_0$, mostre
que $c/c_0 \cong K$ isometricamente, via $\bar x \mapsto
\lim_nx_n$.
\end{exercise}

\begin{exercise}[$\star\star$]\label{exo:b3:banach:12}
(Projeções limitadas e subespaços complementados) Sejam $E$ um
espaço de Banach e $P \colon E \to E$ linear com $P^2 = P$ (uma projeção \omterm{def:b3:galois:algebraic}{algébrica}),
$V = \operatorname{im}P$, $W =
\ker P$.
(a) Suponha $P$ limitada. Mostre que $V$ e $W$ são fechados e
que $E = V \oplus W$, com a decomposição $x = Px + (x -
Px)$.
(b) Reciprocamente, suponha $E = V \oplus W$ com $V,
W$ \emph{ambos} fechados, e
seja $P$ a projeção sobre $V$ ao longo de $W$. Mostre que
$P$ é limitada. \emph{(Gráfico fechado: se $x_n \to x$ e $Px_n \to z$, então
$z \in V$, $x_n - Px_n \to x - z \in
W$, e a unicidade da decomposição identifica $z =
Px$.)}
(c) Deduza a \emph{equivalência}: um subespaço $V$ admite projeção
limitada se, e somente se, é fechado e tem complemento \omterm{def:b3:galois:algebraic}{algébrico}
fechado --- e note (sem demonstração) que existem subespaços fechados
sem essa propriedade ($c_0$ dentro de $\ell^\infty$ é o exemplo clássico):
os espaços de Hilbert, em que $V^\perp$ sempre funciona (\cref{ch:b3:hilbert}), são
a exceção, não a regra.
\end{exercise}

\section{Problema: a dualidade dos espaços
\texorpdfstring{$\ell^p$}%
{lp}}

\begin{problem}[{Problema de fim de semana --- $(\ell^p)' = \ell^q$, reflexividade e a
estranheza de $\ell^\infty$}]
\label{pb:b3:banach:1}
Fixe $1 < p < \infty$ e seja $q$ o expoente conjugado, $\frac1p + \frac1q = 1$. O emparelhamento
em todo o problema é $\langle x,
y\rangle = \sum_n x_ny_n$.

\medskip
\noindent\textbf{Parte I --- Hölder e Minkowski para sequências.}
\begin{enumerate}
  \item (Desigualdade de Young) Para $a, b \geq 0$, mostre que $ab \leq
        \frac{a^p}p + \frac{b^q}q$, usando a
        concavidade de $\log$ ou estudando $t \mapsto \frac{t^p}p + \frac1q - t$.
  \item (Hölder) Deduza: $\abs{\langle x, y\rangle} \leq
        \norm x_p\norm y_q$ para $x \in \ell^p$, $y \in \ell^q$; identifique o caso de
        igualdade.
  \item (Minkowski) Deduza a desigualdade triangular para $\norm\cdot_p$.
        \emph{(Escreva $\abs{x_n + y_n}^p \leq
        \abs{x_n}\,\abs{x_n{+}y_n}^{p-1} +
        \abs{y_n}\,\abs{x_n{+}y_n}^{p-1}$ e aplique Hölder a cada termo.)}
  \item Demonstre que $\ell^p$ é \omterm{def:b3:complete:complete}{completo} e que as sequências
        finitas são densas nele.
\end{enumerate}

\medskip
\noindent\textbf{Parte II --- A dualidade $(\ell^p)' =
\ell^q$.}
\begin{enumerate}[resume]
  \item Para $y \in \ell^q$, mostre que $\Lambda_y(x) = \langle
        x, y\rangle$ define $\Lambda_y \in (\ell^p)'$ com $\norm{\Lambda_y} \leq \norm y_q$ e, testando em
        $x_n = \abs{y_n}^{q-1}\operatorname{sign}(y_n)$ (convenientemente truncado e normalizado), que $\norm{\Lambda_y} = \norm y_q$.
  \item Reciprocamente, dado $\Lambda \in (\ell^p)'$, ponha $y_n =
        \Lambda(e_n)$; mostre que $y \in \ell^q$ com
        $\norm y_q \leq
        \norm\Lambda$ \emph{(teste em truncamentos como na questão 5 e faça o
        comprimento do truncamento crescer)}, e conclua $\Lambda = \Lambda_y$: a
        aplicação $y \mapsto \Lambda_y$ é um isomorfismo isométrico $\ell^q \to (\ell^p)'$.
  \item Deduza que $\ell^p$ é \emph{\omterm{rem:b3:banach:bidual}{reflexivo}} para $1 < p <
        \infty$: compondo
        as duas dualidades, todo elemento de $(\ell^p)''$ provém de
        $\ell^p$; verifique com cuidado que a composta é a
        aplicação canônica $J$.
\end{enumerate}

\medskip
\noindent\textbf{Parte III --- $\ell^1$ e $\ell^\infty$ são animais
diferentes.}
\begin{enumerate}[resume]
  \item Mostre que $\ell^p$ ($1 \leq p < \infty$) e $c_0$ são separáveis, mas
        $\ell^\infty$ não é. \emph{(As não enumeráveis sequências indicadoras
        de subconjuntos de $\N$ estão duas a duas a distância
        $1$.)}
  \item Deduza do~\cref{exo:b3:banach:10} que $(\ell^1)'
        \cong \ell^\infty$, mas $(\ell^\infty)' \not\cong
        \ell^1$: $\ell^1$
        \emph{não} é \omterm{rem:b3:banach:bidual}{reflexivo}. \emph{(Se $(\ell^\infty)'$ fosse $\ell^1$, ele
        seria separável, o que forçaria $\ell^\infty$ separável.)}
  \item (Um limite de Banach, explicitamente) Em $\ell^\infty_\R$, ponha $p(x) = \limsup_n \frac{x_1 + \dots + x_n}{n}$.
        Mostre que $p$ é sublinear e que, no subespaço $c$ das
        sequências convergentes, $\mathrm{LIM}(x) = \lim x$ satisfaz $\mathrm{LIM} \leq p$. Estenda por
        Hahn--Banach a $\mathrm{LIM} \colon \ell^\infty_\R \to \R$ e mostre: $\mathrm{LIM}$ é positivo ($x \geq 0 \Rightarrow
        \mathrm{LIM}(x) \geq 0$),
        invariante por deslocamento ($\mathrm{LIM}(x_2, x_3, \dots) = \mathrm{LIM}(x)$), estende o limite e
        satisfaz $\liminf x \leq
        \mathrm{LIM}(x)\leq \limsup x$.
  \item Mostre que um tal $\mathrm{LIM}$, visto em $(\ell^\infty)'$, \emph{não} é da forma
        $\Lambda_y$ para nenhum $y \in \ell^1$; conclua de novo que $(\ell^\infty)'
        \neq \ell^1$.
        \emph{(Avalie nas sequências unitárias $e_n$ e, depois, na
        sequência constante $1$.)}
  \item Avalie $\mathrm{LIM}$ em $(0,1,0,1,\dots)$ e mostre que não pode existir extensão
        \emph{multiplicativa} do limite que seja invariante por
        deslocamento \emph{(considere $x =
        (0,1,0,1,\dots)$ e $x\cdot Sx$, em que $S$ é o
        deslocamento)}.
\end{enumerate}

\medskip
\noindent\textbf{Parte IV --- Epílogo: por que a reflexividade
importa.}
\begin{enumerate}[resume]
  \item Usando o~\cref{cor:b3:banach:bslimits} e a questão 6, mostre que toda
        sequência limitada de $\ell^p$ ($1 < p <
        \infty$) tem uma subsequência
        $(x^{(k)})$ que converge \emph{fracamente}: $\Lambda(x^{(k)})$ converge para todo
        $\Lambda \in (\ell^p)'$. \emph{(Extração diagonal nas enumeráveis coordenadas;
        identifique o limite fraco em $\ell^p$ usando a limitação
        uniforme das normas e Hölder.)} Mostre com um exemplo
        ($e_n$ em $\ell^1$, contra elementos bem escolhidos de
        $\ell^\infty$) que isso falha em $\ell^1$: a \omterm{def:b3:topology:compact}{compacidade} fraca é um
        privilégio dos \omterm{rem:b3:banach:bidual}{espaços reflexivos}.
\end{enumerate}

\medskip
\noindent\textbf{Parte V --- A \omterm{def:b3:topology:topology}{topologia} fraca em ação, e a surpresa
de Schur.} Escreva $x^{(k)} \rightharpoonup x$ em um espaço normado $E$
(\emph{convergência fraca}) quando $\Lambda(x^{(k)}) \to \Lambda(x)$ para todo $\Lambda \in E'$.
\begin{enumerate}[resume]
  \item Complete o censo: mostre que $(c_0)' \cong \ell^1$ isometricamente, pelo
        esquema das questões 5--6 (o que substitui as sequências de
        teste?). Monte a cadeia $c_0 \to \ell^1 \to \ell^\infty \to \dots$ de duais sucessivos e assinale
        onde a reflexividade falha.
  \item Mostre que toda sequência fracamente convergente de um espaço
        de Banach é limitada: veja os $x^{(k)}$, através do mergulho
        canônico $J$, como funcionais em $E'$ e aplique
        Banach--Steinhaus (\cref{thm:b3:banach:banachsteinhaus}) --- em qual espaço de Banach,
        e por que a \omterm{def:b3:complete:complete}{completude} está disponível ali?
  \item Mostre que, em $\ell^p$ com $1 < p < \infty$: $x^{(k)}
        \rightharpoonup x$ se, e somente se,
        $\sup_k\norm{x^{(k)}}_p <
        \infty$ e $x^{(k)}_n \to x_n$ para toda coordenada $n$ \emph{(uma das
        implicações usa Banach--Steinhaus através do mergulho
        canônico; para a outra, aproxime $y
        \in \ell^q$ por sequências finitas)}.
        Deduza $e_k
        \rightharpoonup 0$ em $\ell^2$, embora $\norm{e_k}_2 =
        1$: os limites fracos
        podem perder massa.
  \item Mostre que a norma é fracamente semicontínua inferiormente:
        $x^{(k)} \rightharpoonup x$ implica $\norm x \leq
        \liminf_k\,\norm{x^{(k)}}$ \emph{(tome um funcional que norme $x$,
        \cref{cor:b3:banach:hbnormed})}.
  \item (Radon--Riesz em $\ell^2$) Mostre que, em $\ell^2$, a
        convergência fraca junto com a convergência das normas implica
        convergência em norma \emph{(expanda $\norm{x^{(k)} -
        x}_2^2$)}. Dê um
        contraexemplo para o mesmo enunciado sem a hipótese sobre as
        normas.
  \item (Schur, etapa 1) Seja $x^{(k)} \rightharpoonup 0$ em $\ell^1$ e suponha, por
        absurdo, $\norm{x^{(k)}}_1 \geq \delta > 0$ ao longo de uma subsequência. Mostre primeiro
        que $x^{(k)}_n \to 0$ para cada $n$ (quais funcionais?), e depois
        construa recursivamente índices $k_1 < k_2 < \cdots$ e inteiros $0 = N_0 < N_1 < N_2 < \cdots$ tais que
        a massa de $x^{(k_j)}$ se concentre no bloco $B_j =
        \intoc{N_{j-1}}{N_j}$:
        \[
        \sum_{n \in B_j}\bigl|x^{(k_j)}_n\bigr| \geq
        \norm{x^{(k_j)}}_1 - \frac\delta{10} .
        \]
  \item (Schur, etapa 2) Defina $y \in \ell^\infty$ por $y_n =
        \operatorname{sign}\bigl(x^{(k_j)}_n\bigr)$ para $n \in
        B_j$. Mostre que
        $\langle x^{(k_j)}, y\rangle \geq
        \norm{x^{(k_j)}}_1 - \frac{2\delta}{10} \geq
        \frac{8\delta}{10}$ e obtenha uma contradição com $x^{(k)} \rightharpoonup 0$. Conclua o
        \emph{teorema de Schur}: em $\ell^1$, as sequências
        fracamente convergentes convergem em norma.
  \item Deduza que $(e_k)$ não tem subsequência fracamente
        convergente em $\ell^1$ (seu único candidato a limite é
        $0$, coordenada a coordenada --- aplique então Schur),
        recuperando a falha de \omterm{def:b3:topology:compact}{compacidade} fraca da questão 13; e
        resolva o aparente paradoxo: em $\ell^1$ as convergências
        fraca e em norma de \emph{sequências} coincidem, e no entanto
        as \emph{\omterm{def:b3:topology:topology}{topologias}} fraca e da norma diferem e os conjuntos
        limitados continuam a não ser fracamente sequencialmente
        \omterm{def:b3:topology:compact}{compactos} --- não há contradição, apenas a falha da
        reflexividade.
  \item (Tabela de síntese) Para $E \in \{c_0,\ \ell^1,\ \ell^p\
        (1{<}p{<}\infty),\ \ell^\infty\}$, tabule: o \omterm{def:b3:banach:operator}{dual}; a
        separabilidade; a reflexividade; se as sequências limitadas
        admitem subsequências fracamente convergentes; e uma
        propriedade característica de cada espaço, justificada em uma
        linha a partir deste problema.
\end{enumerate}

\medskip
\noindent\textbf{Parte VI --- Complementos: pontos mais próximos,
convergência em média e o valor de um limite de Banach.}
\begin{enumerate}[resume]
  \item (Pontos mais próximos: um dividendo da reflexividade) Sejam
        $F$ um subespaço fechado de $\ell^p$ ($1 < p < \infty$) e $x \in \ell^p$.
        Mostre que $d = \operatorname{dist}(x,
        F)$ é \emph{atingida}: extraia de uma sequência
        minimizante uma subsequência fracamente convergente (questão
        13), mantenha o limite fraco dentro de $F$ construindo, via
        Hahn--Banach, um funcional que se anula em $F$ mas não em um
        ponto fora dele, e conclua com a questão 17. Mostre em seguida
        que o privilégio não é universal: em $c_0$, para $\Lambda(x) = \sum_n2^{-n}x_n$,
        demonstre que $\norm\Lambda = 1$ não é atingida na bola unitária, estabeleça
        a fórmula de distância $\operatorname{dist}(x, \ker\Lambda) =
        \abs{\Lambda(x)}$ e deduza que nenhum $x \notin
        \ker\Lambda$ tem ponto
        mais próximo no hiperplano fechado $\ker\Lambda$.
  \item (Banach--Saks em $\ell^2$) Seja $x^{(k)}
        \rightharpoonup 0$ em $\ell^2_{\R}$ com $\norm{x^{(k)}}_2 \leq C$.
        Construa uma subsequência $(y_j)$ com $\abs{\langle y_i, y_j\rangle} \leq
        \frac1j$ para todos
        $i < j$, e deduza
        \[
        \Bigl\lVert\frac{y_1 + \dots + y_m}m\Bigr\rVert_2^2
        \leq \frac{C^2 + 2}m \longrightarrow 0 :
        \]
        após extração, as médias de Cesàro convergem em \emph{norma}.
        Confira em $(e_k)$, cujas médias têm norma $\frac1{\sqrt m}$: a
        convergência fraca, inútil para a própria sequência (questão
        16), torna-se convergência em norma para as médias.
  \item (O valor de um limite de Banach) Sejam $L$ um limite de
        Banach qualquer (questão 10) e $A_mx = \frac1m(x + Sx + \dots
        + S^{m-1}x)$. Mostre que $L(A_mx) = L(x)$ e $\liminf A_mx
        \leq L(x) \leq \limsup A_mx$;
        deduza que \emph{todos} os limites de Banach coincidem nas
        sequências periódicas, com valor a média sobre um período ---
        $\frac13$ em $(1, 0, 0, 1,
        0, 0, \dots)$, coerente com o $\frac12$ da questão 12. Mostre em
        seguida que a coincidência falha em geral: para a sequência em
        blocos $x$ igual a $1$ em $\intoc{3^{j-1}}{3^j}$ para $j$ par e a
        $0$ nos demais índices, mostre que as médias de Cesàro
        oscilam entre $\leq
        \frac13$ e $\geq \frac23$, e construa dois limites de Banach
        $L_\pm$ com $L_-(x) \leq \frac13 < \frac23
        \leq L_+(x)$ \emph{(estenda a partir de $c \oplus \R x$ com os
        valores admissíveis extremos $\pm$: verifique que $\Lambda(y + tx) = \lim y + t\,p(x)$ é
        dominado pelo $p$ sublinear da questão 10)}.
\end{enumerate}
\end{problem}
